\chapter{Stroomkringen en veiligheid}\label{ch:g6:circuits-and-safety}

Drie jaar lussen, \omterm{def:g3:first-electric-circuit:bulb}{lampjes} en \omterm{def:g5:series-parallel-circuits:parallel}{takken} hebben van je een bouwer van
schakelingen gemaakt. Dit jaar word je iets zeldzamers: een
\emph{schrijver} van schakelingen. Elektriciens in elk land lezen en
tekenen dezelfde woordeloze taal --- en haar grammatica, plus een
hoofdstuk vol duur betaalde veiligheidswijsheid, is wat er staat tussen
slim bedraden en verbrande vingers.

\section{De woordeloze taal}

\begin{definition}[Schakelschema]\label{def:g6:circuits-and-safety:diagram}
Een \emph{schakelschema}\index{schakelschema} is de standaardtekening
van een schakeling: elk apparaat wordt weergegeven door zijn
afgesproken \emph{symbool}\index{symbool}, en de draden door rechte
lijnen met rechte hoeken. Een schema laat zien hoe de schakeling
\emph{verbonden} is --- wie wie voedt, wat waar vertakt --- en negeert
met opzet hoe de echte draden over de tafel kronkelen.
\end{definition}

\begin{omfigure}
\begin{tikzpicture}[scale=0.8]
  \draw (0.4,3.4) to[battery1] (2.4,3.4);
  \node[right, font=\small] at (2.7,3.4)
    {batterij (lang streepje: $+$)};
  \draw (0.4,2.2) to[lamp] (2.4,2.2);
  \node[right, font=\small] at (2.7,2.2) {lamp};
  \draw (0.4,1.0) to[spst] (2.4,1.0);
  \node[right, font=\small] at (2.7,1.0) {schakelaar, open};
  \draw (7.4,3.4) to[cspst] (9.4,3.4);
  \node[right, font=\small] at (9.7,3.4) {schakelaar, dicht};
  \draw (7.4,2.2) -- (8.05,2.2);
  \draw[thick] (8.4,2.2) circle (0.35);
  \node[font=\small] at (8.4,2.2) {M};
  \draw (8.75,2.2) -- (9.4,2.2);
  \node[right, font=\small] at (9.7,2.2) {motor};
  \draw (7.4,1.0) -- (8.05,1.0);
  \draw[thick] (8.4,1.0) circle (0.35);
  \draw[thick] (8.16,0.76) -- (8.64,1.24);
  \draw (8.75,1.0) -- (9.4,1.0);
  \node[right, font=\small] at (9.7,1.0) {zoemer};
\end{tikzpicture}

{\small Het alfabet: zes symbolen dragen bijna elke schakeling van dit
boek. Dezelfde tekens worden van werkplaats tot werkplaats over de
hele wereld gelezen.}
\end{omfigure}

\begin{method}[Van tafel naar papier, en terug]\label{met:g6:circuits-and-safety:translate}
Om een gebouwde schakeling te \emph{tekenen}:
\begin{enumerate}
  \item maak een lijst van de apparaten en kies van elk het
        \omterm{def:g6:circuits-and-safety:diagram}{symbool};
  \item volg de echte lus met je vinger vanaf de $+$ van de
        \omterm{def:g3:first-electric-circuit:battery}{batterij}, en noteer de volgorde van de apparaten en elke
        vertakking;
  \item teken het opnieuw als een nette rechthoek: symbolen uit
        elkaar, draden recht, hoeken haaks --- dezelfde
        verbindingen, geen slangen.
\end{enumerate}
Om vanuit een schema te \emph{bouwen}, laat je de vertaling
achterstevoren lopen --- en controleer daarna, zoals altijd: elke lus
die je vanaf $+$ volgt moet via minstens één apparaat bij $-$
aankomen.
\end{method}

\begin{example}[Dezelfde schakeling, twee portretten]\label{ex:g6:circuits-and-safety:portraits}
Op tafel: een \omterm{def:g3:first-electric-circuit:battery}{batterij}, een kluwen krokodillenklemdraden die over
een etui lussen, een lamp die tegen een boek leunt, een \omterm{ex:g3:first-electric-circuit:switch}{schakelaar}
die erbij bungelt. Op papier: een rustige rechthoek --- \omterm{def:g3:first-electric-circuit:battery}{batterij}
links, \omterm{ex:g3:first-electric-circuit:switch}{schakelaar} bovenaan, lamp rechts. De kluwen en de rechthoek
zijn \emph{dezelfde schakeling}: elke verbinding identiek. Om te
beslissen of twee schakelingen dezelfde zijn, stel je maar één vraag:
wie is met wie verbonden?
\end{example}

\begin{example}[Motoren en zoemers doen mee]\label{ex:g6:circuits-and-safety:motor}
Twee nieuwkomers in je kist gedragen zich in een schema precies als
lampen: de \emph{motor}\index{motor} (\omterm{def:g6:circuits-and-safety:diagram}{symbool}: M in een cirkel)
draait zodra er stroom doorheen loopt --- ventilatoren, speelgoedauto's,
boormachines; de \emph{zoemer}\index{zoemer} (een aangeslagen cirkel)
klinkt zodra hij gevoed wordt --- deurbellen, alarmen. De regels van
serie en \omterm{def:g5:series-parallel-circuits:parallel}{parallel} gelden onveranderd voor hen: een \omterm{ex:g3:first-electric-circuit:switch}{schakelaar}
\omterm{def:g4:series-circuits:series}{in serie} commandeert ze; parallelle \omterm{def:g5:series-parallel-circuits:parallel}{takken} maken ze
onafhankelijk.
\end{example}

\section{Oude regels in de nieuwe taal}

\begin{example}[Een schema lezen als een vakman]\label{ex:g6:circuits-and-safety:reading}
Ondervraag elk schema met de vingerwandeling uit het hoofdstuk over
\omterm{def:g5:series-parallel-circuits:parallel}{parallelschakelingen}: loop vanaf $+$ elke weg naar $-$. Alle apparaten
op één weg: serie --- één onderbreking legt alles stil, de volgorde is
vrij, de felheid wordt gedeeld. Wegen die bij knooppuntstippen
vertakken: \omterm{def:g5:series-parallel-circuits:parallel}{parallel} --- \omterm{def:g5:series-parallel-circuits:parallel}{takken} onafhankelijk, elk op volle felheid, en
de plaatsing van de \omterm{ex:g3:first-electric-circuit:switch}{schakelaar} is nu een echte beslissing. De taal
veranderde; de wetten niet.
\end{example}

\begin{omfigure}
\begin{tikzpicture}[scale=0.8]
  \draw (0,0) to[battery1] (0,3.2)
    to[cspst] (2.6,3.2);
  \draw (2.6,3.2) to[short] (2.6,2.4)
    to[lamp] (5.8,2.4) to[short] (5.8,3.2);
  \draw (2.6,3.2) to[short] (2.6,4.0) -- (3.4,4.0);
  \draw[thick] (3.75,4.0) circle (0.35);
  \node[font=\small] at (3.75,4.0) {M};
  \draw (4.1,4.0) -- (5.8,4.0) to[short] (5.8,3.2);
  \draw (5.8,3.2) to[short] (7.4,3.2) to[short] (7.4,0)
    to[short] (0,0);
  \fill (2.6,3.2) circle (2pt);
  \fill (5.8,3.2) circle (2pt);
\end{tikzpicture}

{\small Lees het voordat je het bouwt: één \omterm{ex:g3:first-electric-circuit:switch}{schakelaar} op de
gemeenschappelijke weg, en dan een vertakking --- lamp op de ene
\omterm{def:g5:series-parallel-circuits:parallel}{tak}, ventilatormotor op de andere. De \omterm{ex:g3:first-electric-circuit:switch}{schakelaar} commandeert
allebei; de \omterm{def:g5:series-parallel-circuits:parallel}{takken} negeren elkaars problemen.}
\end{omfigure}

\section{Het hoofdstuk van de littekens}

De veiligheidsregels van de elektriciteit zijn betaald met branden en
erger. Hier zijn ze, met hun redenen --- een vakman kent allebei.

\begin{definition}[Kortsluiting]\label{def:g6:circuits-and-safety:short}
Van \emph{kortsluiting}\index{kortsluiting} is sprake wanneer een
blote geleidende weg de twee kanten van een \omterm{def:g3:first-electric-circuit:battery}{batterij} (of van een
stopcontact) \emph{rechtstreeks} verbindt, met alle apparaten
overgeslagen. Met niets nuttigs in de weg stormt de stroom
ongeremd --- draden verhitten, isolatie \omterm{def:g2:ice-water-steam:meltfreeze}{smelt}, \omterm{def:g3:first-electric-circuit:battery}{batterijen} zijn
verwoest, en netbedrading kan vlam vatten. Het schakelingenhoofdstuk
van volgend jaar bestudeert de schurk van dichtbij; leer dit jaar zijn
gezicht kennen en houd hem uit je bouwsels.
\end{definition}

\begin{remark}[De lijfwachten van het huis]\label{rem:g6:circuits-and-safety:guards}
Kijk thuis in de meterkast: rijen kleine \omterm{ex:g3:first-electric-circuit:switch}{schakelaars} ---
\emph{stroomonderbrekers} --- één per \omterm{def:g5:series-parallel-circuits:parallel}{tak} van de grote
\omterm{def:g5:series-parallel-circuits:parallel}{parallelschakeling} van het huis. Elk houdt de wacht over zijn
\omterm{def:g5:series-parallel-circuits:parallel}{tak}: laat een storing de stroom op hol slaan --- een
\omterm{def:g6:circuits-and-safety:short}{kortsluiting}, een overbelast stopcontact --- dan klapt de wachter
zichzelf in een oogwenk uit en snijdt de \omterm{def:g5:series-parallel-circuits:parallel}{tak} af voordat draden
oververhit raken. Wanneer een onderbreker ``eruit vliegt'', misdraagt
hij zich niet: hij heeft zojuist zijn enige taak gedaan. Zoek de
storing voordat je hem terugzet.
\end{remark}

\begin{remark}[De regels, met redenen]\label{rem:g6:circuits-and-safety:rules}
\begin{enumerate}
  \item \emph{Alleen batterijproeven; het stopcontact nooit.} Het
        net duwt honderden keren harder --- door vochtige huid kan
        het een hart stilleggen.
  \item \emph{Geen blote draad dwars over een \omterm{def:g3:first-electric-circuit:battery}{batterij}.} Dat is
        \omterm{def:g6:circuits-and-safety:short}{kortsluiting}: verwoeste cellen, hete draad, verbrande
        vingers.
  \item \emph{Water en netstroom komen nooit binnen elkaars bereik:}
        geen apparaten bij het bad, geen natte handen aan stekkers
        --- vochtige huid is een \omterm{def:g4:conductors-insulators:conductor}{geleider}.
  \item \emph{Steek nooit iets in een stopcontact.} De wachter in de
        meterkast is snel, maar je mag er nooit op wedden.
  \item \emph{Beschadigde snoeren gaan naar de vuilnisbak,} niet naar
        de la: gebarsten isolatie is een hek met een gat.
\end{enumerate}
\end{remark}

\section{Oefeningen}

\begin{exercise}[$\star$]\label{exo:g6:circuits-and-safety:1}
Wat laat een \omterm{def:g6:circuits-and-safety:diagram}{schakelschema} getrouw zien, en wat negeert het met
opzet?
\end{exercise}

\begin{exercise}[$\star$]\label{exo:g6:circuits-and-safety:2}
Teken uit je hoofd de symbolen: \omterm{def:g3:first-electric-circuit:battery}{batterij}; \ lamp; \ open
\omterm{ex:g3:first-electric-circuit:switch}{schakelaar}; \ gesloten \omterm{ex:g3:first-electric-circuit:switch}{schakelaar}; \ \omterm{ex:g6:circuits-and-safety:motor}{motor}; \ \omterm{ex:g6:circuits-and-safety:motor}{zoemer}. Welk teken
verraadt de $+$-kant van de \omterm{def:g3:first-electric-circuit:battery}{batterij}?
\end{exercise}

\begin{exercise}[$\star$]\label{exo:g6:circuits-and-safety:3}
Teken het schema van een deurbel: \omterm{def:g3:first-electric-circuit:battery}{batterij}, drukschakelaar en
\omterm{ex:g6:circuits-and-safety:motor}{zoemer} in één lus. Waarom stopt de bel zodra de vinger loslaat?
\end{exercise}

\begin{exercise}[$\star$]\label{exo:g6:circuits-and-safety:4}
Twee schakelingen worden uit dezelfde doos onderdelen gebouwd: zijn ze
``dezelfde schakeling''? Welke ene vraag beslist dat?
\end{exercise}

\begin{exercise}[$\star$]\label{exo:g6:circuits-and-safety:5}
Wat gebeurt er in het lamp-en-motorschema van het hoofdstuk met de
ventilator als de lamp doorbrandt? En met allebei als de
\omterm{ex:g3:first-electric-circuit:switch}{schakelaar} opengaat? Welke oude regels geven het antwoord?
\end{exercise}

\begin{exercise}[$\star$]\label{exo:g6:circuits-and-safety:6}
Wat is \omterm{def:g6:circuits-and-safety:short}{kortsluiting}, en waarom verhit ze draden? Noem twee van haar
slachtoffers.
\end{exercise}

\begin{exercise}[$\star$]\label{exo:g6:circuits-and-safety:7}
Een stroomonderbreker in de meterkast is eruit gevlogen. Wat heeft hij
zojuist voor het huis gedaan? Wat moet er gebeuren voordat je hem
terugzet?
\end{exercise}

\begin{exercise}[$\star$]\label{exo:g6:circuits-and-safety:8}
Geef de reden achter elke regel: geen apparaten bij het bad; \
beschadigde snoeren naar de vuilnisbak; \ batterijspelletjes ja,
stopcontact nooit.
\end{exercise}

\begin{exercise}[$\star\star$]\label{exo:g6:circuits-and-safety:9}
Vertaal naar een schema (in woorden of in een tekening): een
\omterm{def:g3:first-electric-circuit:battery}{batterij}; vanaf haar $+$ een draad naar een gesloten
\omterm{ex:g3:first-electric-circuit:switch}{schakelaar}; vanaf de \omterm{ex:g3:first-electric-circuit:switch}{schakelaar} een vertakking --- één
\omterm{def:g5:series-parallel-circuits:parallel}{tak} door een \omterm{ex:g6:circuits-and-safety:motor}{zoemer}, één door een lamp; de \omterm{def:g5:series-parallel-circuits:parallel}{takken} komen
weer samen en keren terug naar $-$. Voorspel wat je hoort en ziet, en
wat er verandert als alleen de \omterm{def:g5:series-parallel-circuits:parallel}{tak} van de \omterm{ex:g6:circuits-and-safety:motor}{zoemer} een eigen extra
\omterm{ex:g3:first-electric-circuit:switch}{schakelaar} krijgt die openstaat.
\end{exercise}

\begin{exercise}[$\star\star$]\label{exo:g6:circuits-and-safety:10}
Een speelgoedventilator is zo bedraad: batterij-$+$ naar de
\omterm{ex:g6:circuits-and-safety:motor}{motor}, \omterm{ex:g6:circuits-and-safety:motor}{motor} naar batterij-$-$, en --- ``voor de stevigheid''
--- één extra blote draad rechtstreeks van $+$ naar $-$. De
ventilator draait nauwelijks en de extra draad wordt heet.
Diagnosticeer met \cref{def:g6:circuits-and-safety:short}: welke weg
verkiest de stroom, en waarom is het bouwsel gevaarlijk?
\end{exercise}

\begin{exercise}[$\star\star$]\label{exo:g6:circuits-and-safety:11}
De werkplaatstruc van opa: voordat hij een lamparmatuur aanraakt, zet
hij de wandschakelaar \emph{én} de stroomonderbreker van de hele
\omterm{def:g5:series-parallel-circuits:parallel}{tak} in de meterkast uit. Leg met serieredeneren uit wat elke
onderbreking bereikt --- en waarom de voorzichtige man ze allebei
maakt.
\end{exercise}

\begin{exercise}[$\star\star\star$]\label{exo:g6:circuits-and-safety:12}
Ontwerp de bedrading van een poppentheater: één hoofdschakelaar die
alles verduistert; twee toneellampen die onafhankelijk van elkaar
kapot moeten kunnen gaan; een gordijnmotor en een
voorstellingszoemer, elk met een eigen bedieningsschakelaar.
Beschrijf (of teken) het schema --- welke apparaten op de
gemeenschappelijke weg rijden, welke op \omterm{def:g5:series-parallel-circuits:parallel}{takken}, waar elke
\omterm{ex:g3:first-electric-circuit:switch}{schakelaar} zit --- en toets je ontwerp aan elke regel van serie en
\omterm{def:g5:series-parallel-circuits:parallel}{parallel} die je bezit.
\end{exercise}

\section{Probleem: het tuinhuisje bedraden}

\begin{problem}[{Weekendprobleem --- een gezin bedraadt zijn
tuinhuisje; schema's vóór schroevendraaiers; de controlelijst van de
inspecteur}]\label{pb:g6:circuits-and-safety:1}
Het gezin plant het kleine netwerk op \omterm{def:g3:first-electric-circuit:battery}{batterijen} van het tuinhuisje
eerst op papier --- zoals vakmensen doen.

\medskip\noindent\textbf{Deel I --- Het plan op papier.}
Het verlanglijstje: één plafondlamp die bij de deur wordt geschakeld;
een werkbanklamp met een eigen \omterm{ex:g3:first-electric-circuit:switch}{schakelaar}; een kleine
ventilatieventilator (\omterm{ex:g6:circuits-and-safety:motor}{motor}) met een eigen \omterm{ex:g3:first-electric-circuit:switch}{schakelaar}.
\begin{enumerate}
  \item Waarom moeten de drie apparaten op parallelle \omterm{def:g5:series-parallel-circuits:parallel}{takken} rijden
        in plaats van in één serielus? Noem twee gemakken die
        \omterm{def:g4:series-circuits:series}{in serie} zouden sneuvelen.
  \item Elke \omterm{def:g5:series-parallel-circuits:parallel}{tak} draagt zijn eigen \omterm{ex:g3:first-electric-circuit:switch}{schakelaar}. Wat commandeert
        elke \omterm{ex:g3:first-electric-circuit:switch}{schakelaar}, en wat commandeert geen van hen?
  \item Waar moet een hoofdschakelaar zitten om het hele huisje in
        één keer donker en stil te maken?
  \item Teken (of beschrijf nauwkeurig) het volledige schema:
        \omterm{def:g3:first-electric-circuit:battery}{batterij}, hoofdschakelaar, drie \omterm{def:g5:series-parallel-circuits:parallel}{takken}, vier
        \omterm{ex:g3:first-electric-circuit:switch}{schakelaars} in totaal.
\end{enumerate}

\medskip\noindent\textbf{Deel II --- De bouw, geïnspecteerd.}
\begin{enumerate}[resume]
  \item Gebouwd en ingeschakeld schijnt de plafondlamp --- maar
        zwak, en de werkbanklamp schijnt alleen als de
        \emph{ventilator}\omterm{ex:g3:first-electric-circuit:switch}{schakelaar} dicht is. De vingerwandeling
        vindt de werkbanklamp en de ventilator op één \omterm{def:g5:series-parallel-circuits:parallel}{tak}, achter
        elkaar. Noem hun opstelling, en verklaar beide verschijnselen
        met de oude regels.
  \item Geef de reparatie van één draad die het plan herstelt.
  \item Tijdens het opruimen ligt een schroevendraaier even dwars
        over de twee blote polen van de \omterm{def:g3:first-electric-circuit:battery}{batterij} en wordt warm.
        Noem het verschijnsel, en zeg wat het gezin daarna over de
        \omterm{def:g3:first-electric-circuit:battery}{batterij} moet controleren.
  \item De inspecterende ouder voegt een regel toe: ``elke las
        omwikkeld, nergens bloot koper.'' Tegen welke twee gevaren
        waakt die ene regel?
\end{enumerate}

\medskip\noindent\textbf{Deel III --- De nacht van de storm.}
\begin{enumerate}[resume]
  \item Een regenachtige nacht: het dak van het huisje lekt op de
        werkbank. Waarom mag het gezin de bedrading niet met natte
        handen aanraken, zelfs in een huisje op \omterm{def:g3:first-electric-circuit:battery}{batterijen} --- en
        waarom zou diezelfde aanraking in het huis op netstroom veel
        ernstiger zijn?
  \item De lampen van het huis flikkeren en er vliegt een
        stroomonderbreker uit. Wat heeft de onderbreker opgemerkt, en
        wat is de verstandige volgorde van handelingen?
  \item Ochtend: het gezin overweegt een deurbel toe te voegen ---
        \omterm{ex:g6:circuits-and-safety:motor}{zoemer} bij het huisje, drukschakelaar bij het huis, een
        lange dubbele draad ertussen. \omterm{def:g4:series-circuits:series}{In serie} of \omterm{def:g5:series-parallel-circuits:parallel}{parallel} met het
        netwerk van het huisje? Waar komt zijn stroom vandaan? Schets
        de lus.
  \item Schrijf het veiligheidshandvest van het gezin voor het
        huisje in drie regels, elk één regel: stopcontacten, water,
        blote draden.
\end{enumerate}
\end{problem}
